\section{Introduction}

Buồn ngủ (\textit{drowsiness}) là quá trình suy giảm dần mức độ tỉnh táo (\textit{declining vigilance}), đi kèm với giảm khả năng duy trì chú ý và phản ứng với kích thích bên ngoài \cite{peker2026,el2026}. Photoplethysmography (PPG) phù hợp để theo dõi quá trình này nhờ tính không xâm lấn, chi phí thấp và khả năng tích hợp vào thiết bị đeo. Phần lớn nghiên cứu PPG về drowsiness khai thác các đặc trưng sinh lý, hình thái và miền thời gian--tần số, bao gồm heart-rate variability (HRV), chủ yếu để phục vụ detection/classification \cite{amidei2021,cao2023,alarnaout2025,yu2024}. Bản thân waveform PPG còn chứa tổ chức theo thời gian và các đặc tính động lực học phi tuyến \cite{de2020,sviridova2015}. Trọng tâm phân loại của các cách tiếp cận hiện có đặt ra nhu cầu đặc trưng trực tiếp tổ chức động lực học của waveform theo trạng thái.

Nhóm nghiên cứu nonlinear PPG của Sviridova và cộng sự cung cấp ba cơ sở cho hướng tiếp cận này: tín hiệu thể hiện tổ chức động lực học phi tuyến \cite{sviridova2015}; tính dao động gần chu kỳ thúc đẩy pseudoperiodic surrogate (PPS) testing để kiểm tra tổ chức vượt quá một noisy pseudoperiodic null cụ thể \cite{sviridova2018}; và thông tin động lực học hữu ích có thể được khảo sát từ bản ghi tương đối ngắn \cite{sviridova2022}. Tuy nhiên, cách finite-horizon forecastability, recurrence organization và local trajectory divergence cùng thay đổi qua chuyển tiếp Awake--Drowsy vẫn chưa được đặc trưng bằng một mô tả tích hợp. Khoảng trống này đòi hỏi kết hợp kiểm tra tổ chức vượt quá null đã chọn với đánh giá những thuộc tính động lực học bổ sung nhau.

Ba góc nhìn gồm khả năng dự báo trong khoảng thời gian hữu hạn (\textit{finite-horizon forecastability}), được đánh giá bằng Simplex Projection; tổ chức tái diễn (\textit{recurrence organization}), được mô tả bằng Recurrence Quantification Analysis (RQA); và sự phân kỳ cục bộ của các quỹ đạo lân cận (\textit{local trajectory divergence}), được ước lượng bằng Rosenstein Largest Lyapunov Exponent (LLE). Vì chúng đặc trưng các khía cạnh khác nhau, phân tích phối hợp giúp tránh diễn giải động lực học theo một trục tăng--giảm đơn giản của chaos hay complexity. Cửa sổ ngắn giúp phân giải quá trình suy giảm tỉnh táo theo thời gian, nhưng làm giảm số điểm trong không gian trạng thái, cấu trúc recurrence và cặp quỹ đạo, từ đó tăng độ bất định của các ước lượng. Độ tin cậy cần được xem xét theo từng thuộc tính và độ dài dữ liệu \cite{sviridova2022}. Vì vậy, so sánh trạng thái cần đi kèm kiểm tra độ bền theo độ dài cửa sổ và độ nhạy với các lựa chọn reconstruction/estimation chính.

Nghiên cứu tập trung vào hai câu hỏi: \textbf{RQ1:} Liệu PPG cửa sổ ngắn ở Awake và Drowsy có chứa tổ chức thời gian/động lực học vượt quá noisy pseudoperiodic null đã kiểm tra hay không? \textbf{RQ2:} Drowsiness liên quan như thế nào đến các thay đổi của finite-horizon forecastability, recurrence organization và local trajectory divergence? RQ1 được đánh giá bằng PPS testing; RQ2 từ các khác biệt Awake--Drowsy tổng hợp ở cấp recording session. Độ nhạy với repeated-session dependence chưa biết được xem xét bổ sung, cùng các kiểm tra theo độ dài cửa sổ, định nghĩa nhãn và thiết lập phân tích. Đóng góp chính là đặc trưng hóa sự tái tổ chức phối hợp của PPG dynamics qua chuyển tiếp Awake--Drowsy trên ba thuộc tính bổ sung nhau, đồng thời xác định mức độ bền của các kết quả trong phạm vi độ dài dữ liệu và các lựa chọn phân tích đã kiểm tra.
